
\documentclass{article}
\usepackage{colortbl}
\usepackage{makecell}
\usepackage{multirow}
\usepackage{supertabular}

\begin{document}

\newcounter{utterance}

\centering \large Interaction Transcript for game `cladder', experiment `full\_v1.5\_default', episode 8398 with Qwen3.6{-}27B.
\vspace{24pt}

{ \footnotesize  \setcounter{utterance}{1}
\setlength{\tabcolsep}{0pt}
\begin{supertabular}{c@{$\;$}|p{.15\linewidth}@{}p{.15\linewidth}p{.15\linewidth}p{.15\linewidth}p{.15\linewidth}p{.15\linewidth}}
    \# & \multicolumn{2}{c}{Player} && \multicolumn{2}{c}{Game Master} \\
    \hline

    \theutterance \stepcounter{utterance}  
    & & & \multicolumn{4}{p{0.6\linewidth}}{
        \cellcolor[rgb]{0.9,0.9,0.9}{
            \makecell[{{p{\linewidth}}}]{
                \texttt{\tiny{[P1$\langle$GM]}}
                \texttt{You are an expert in causal inference. The following question is not a typical commonsense query, but rather a meticulously designed question created by a professor specializing in causal inference, intended to assess the students' mastery of the course content.} \\
\texttt{Imagine a self{-}contained, hypothetical world with only the following conditions, and without any unmentioned factors or causal relationships: Jyka has a direct effect on hwax. Hwax has a direct effect on lirg. The overall probability of jyka is 14\%. The probability of not jyka and lirg is 72\%. The probability of jyka and lirg is 12\%. Is the chance of lirg smaller when observing jyka?} \\
\texttt{Start your answer with 'Yes' or 'No', followed by additional reasoning or evidence to support your explanation.} \\
            }
        }
    }
     \\ \\

    \theutterance \stepcounter{utterance}  
    & \multicolumn{4}{p{0.6\linewidth}}{
        \cellcolor[rgb]{0.561,0.737,0.561}{
            \makecell[{{p{\linewidth}}}]{
                \texttt{\tiny{[P1$\rangle$GM]}}
                \texttt{Yes} \\
\\ 
\texttt{To determine whether the chance of `lirg` is smaller when observing `jyka`, we need to compare the conditional probability \$P(\textbackslash{}text\{lirg\} | \textbackslash{}text\{jyka\})\$ with the marginal probability \$P(\textbackslash{}text\{lirg\})\$.} \\
\\ 
\texttt{**Step 1: Calculate the marginal probability of `lirg`, \$P(\textbackslash{}text\{lirg\})\$.**} \\
\texttt{We are given:} \\
\texttt{{-} \$P(\textbackslash{}neg \textbackslash{}text\{jyka\} \textbackslash{}cap \textbackslash{}text\{lirg\}) = 0.72\$} \\
\texttt{{-} \$P(\textbackslash{}text\{jyka\} \textbackslash{}cap \textbackslash{}text\{lirg\}) = 0.12\$} \\
\\ 
\texttt{Since the events \$(\textbackslash{}neg \textbackslash{}text\{jyka\} \textbackslash{}cap \textbackslash{}text\{lirg\})\$ and \$(\textbackslash{}text\{jyka\} \textbackslash{}cap \textbackslash{}text\{lirg\})\$ are mutually exclusive and exhaustive with respect to `lirg` (i.e., `lirg` can only occur with either `jyka` or not `jyka`), we can sum these probabilities to find the total probability of `lirg`:} \\
\texttt{\$\$P(\textbackslash{}text\{lirg\}) = P(\textbackslash{}neg \textbackslash{}text\{jyka\} \textbackslash{}cap \textbackslash{}text\{lirg\}) + P(\textbackslash{}text\{jyka\} \textbackslash{}cap \textbackslash{}text\{lirg\})\$\$} \\
\texttt{\$\$P(\textbackslash{}text\{lirg\}) = 0.72 + 0.12 = 0.84\$\$} \\
\\ 
\texttt{**} \\
            }
        }
    }
    & & \\ \\

    \theutterance \stepcounter{utterance}  
    & & & \multicolumn{2}{p{0.3\linewidth}}{
        \cellcolor[rgb]{0.9,0.9,0.9}{
            \makecell[{{p{\linewidth}}}]{
                \texttt{\tiny{[GM$|$GM]}}
                \texttt{Yes} \\
            }
        }
    }
    & & \\ \\

    \theutterance \stepcounter{utterance}  
    & & & \multicolumn{2}{p{0.3\linewidth}}{
        \cellcolor[rgb]{0.9,0.9,0.9}{
            \makecell[{{p{\linewidth}}}]{
                \texttt{\tiny{[GM$|$GM]}}
                \texttt{no} \\
            }
        }
    }
    & & \\ \\

    \theutterance \stepcounter{utterance}  
    & & & \multicolumn{2}{p{0.3\linewidth}}{
        \cellcolor[rgb]{0.9,0.9,0.9}{
            \makecell[{{p{\linewidth}}}]{
                \texttt{\tiny{[GM$|$GM]}}
                \texttt{game\_result = LOSE} \\
            }
        }
    }
    & & \\ \\

\end{supertabular}
}

\end{document}
